% Caractéristique tension-intensité d'une pile (exercice 2, chapitre 14)
\begin{tikzpicture}
\begin{axis}[otMesures,otLecture,otAxesBout,width=96mm,height=56mm,xmin=0,xmax=3.4,ymin=0,ymax=5.4,
  xtick={0,0.5,...,3.0},ytick={0,1,...,5},
  xticklabel style={/pgf/number format/fixed,/pgf/number format/precision=1},
  xlabel={$I$ (A)},ylabel={$U$ (V)}]
\addplot[otNavy,line width=1.2pt,domain=0:3.0] {4.5-1.5*x};
\addplot[only marks,mark=*,mark size=1.4pt,otOrange!85!black] coordinates {(0,4.5) (0.4,3.9) (0.8,3.3) (1.2,2.7) (1.6,2.1) (2.0,1.5)};
\solution{
  \draw[otGreenDark,line width=0.6pt,dash pattern=on 2pt off 1.5pt] (axis cs:0.4,0) -- (axis cs:0.4,3.9) -- (axis cs:0,3.9);
  \node[pcLegende,anchor=south west,text=otGreenDark] at (axis cs:0.42,3.92) {\qty{3,9}{V} à \qty{0,40}{A}};
  \node[pcLegende,anchor=south west,text=otGreenDark] at (axis cs:0.02,4.52) {$E=\qty{4,5}{V}$};
  \fill[otGreenDark] (axis cs:3.0,0) circle (1.5pt);
  \node[pcLegende,anchor=south west,text=otGreenDark,fill=white,inner sep=1pt] at (axis cs:3.03,0.12) {$I_\mathrm{cc}=\qty{3,0}{A}$};
  \draw[otGreenDark,line width=0.6pt] (axis cs:1.6,2.1) -- (axis cs:2.6,2.1) -- (axis cs:2.6,0.6);
  \node[pcLegende,anchor=west,text=otGreenDark,align=left] at (axis cs:2.62,1.35) {$\dfrac{-1{,}5}{1{,}0}$ : $r=\qty{1,5}{\ohm}$};}
\end{axis}
\end{tikzpicture}
\par\smallskip\legende{Points de mesure et droite modèle de la caractéristique d'une pile.}
